        \section{供应链门禁结论：可以覆盖，但不能无条件上调估值}
        独立产业链研究门禁给出的状态是有条件通过。18只核心标的均已建立公司卡片和关系行，但具名晶圆厂客户、长期订单、产品价格、分产品收入和毛利率拆分仍然不完整。这个状态决定了正文的写法：可以把WF6定义为钨、含氟化学和电子特气的交汇点，也可以把中船特气、中巨芯-U、昊华科技等放入直接验证池；但不能把“下游HBM/3D NAND需求存在”直接写成“上游公司订单和利润已经兑现”。

        对估值最关键的分歧，是利润池兑现速度不同。钨资源和制冷剂的当期财报验证相对直接，能进入基本面锚；电子特气的产业地位更高，但若缺少客户、订单、价格和收入占比，只能进入市场情绪锚或事件驱动锚。换句话说，产业链分析不是给所有标的加分，而是决定哪些标的可以拿到盈利信用，哪些只能拿到观察信用。

        \begin{exhibitbox}[供应链门禁和利润池分层]
        \scriptsize
\resizebox{\exhibitwidth}{!}{
\begin{tabular}{llll}
\toprule
\textbf{利润池} & \textbf{当期兑现} & \textbf{关键证据} & \textbf{估值处理} \\
\midrule
钨资源/冶炼/加工 & 较强 & 钨价、出口管制、APT/钨粉价差、库存成本 & 周期/资源估值可用；半导体外延折扣处理 \\
制冷剂/氟化工 & 较强 & 配额、价差、装置负荷、电子材料认证 & 现金流和电子材料期权分开定价 \\
WF6/NF3/电子特气 & 分化 & 产能、纯度、认证、订单、价格、收入占比 & 无订单/收入占比时仅给事件驱动信用 \\
下游晶圆/设备需求 & 需求验证 & 晶圆厂扩产、CVD/刻蚀设备、HBM/3D NAND & 只作为需求锚，不直接确认上游收入 \\
\bottomrule
\end{tabular}
}
\end{exhibitbox}

\begin{exhibitbox}[政策/出口管制官方证据]
\scriptsize
\begin{tabularx}{\exhibitboxwidth}{L{1.5cm} X X X X}
\toprule
\textbf{政策层} & \textbf{覆盖物项} & \textbf{市场意义} & \textbf{估值用途} & \textbf{证据边界} \\
\midrule
实施时点与法律框架 & 钨、碲、铋、钼、铟等相关物项及部分技术资料 & 把钨从普通周期金属提升为战略资源和出口许可变量 & 支持钨资源安全溢价和出口许可监测，不直接填入公司销售价格或利润率 & 官方政策证明供给约束和战略属性，不证明公司实际成交价、销量、库存成本或毛利 \\
钨相关材料 & APT、氧化钨、碳化钨等钨中间品 & 直接影响APT/钨粉/碳化钨链条的出口许可预期和价格弹性 & 可进入钨链供给约束和资源/加工价差讨论 & 管制范围不等于单家公司可提价，仍需公司销售口径、库存成本和客户结构 \\
固态钨与高密度合金 & 固态钨、高钨含量合金、钨镍铁/钨镍铜合金 & 连接军工、高端制造、高温/高密度材料和出口合规预期 & 支持高端钨材和战略材料外延，但需要公司产品结构验证 & 高端材料政策属性不能替代半导体高纯钨材客户、订单和收入占比 \\
技术与资料 & 生产相关技术、资料、工艺规范、工艺参数、加工程序 & 说明管制不只覆盖商品，也覆盖工艺能力和技术外溢 & 支持产业链安全和国产化能力壁垒叙事 & 技术管制不能替代公司研发成果、客户认证或批量供货证明 \\
政策目的与许可边界 & 出口许可和合规审查框架 & 说明政策不是绝对禁运，市场应跟踪许可与需求，而不是机械外推短缺 & 用于风险折扣和供给约束敏感性，不用于无条件上修倍数 & 许可制不是公司订单，不证明销量、价格或利润已经兑现 \\
\bottomrule
\end{tabularx}
\sourcenote{商务部政策摘要和新闻发言人答问归档整理；政策管制证据只证明出口许可和资源安全属性，不替代公司实际销售价格、订单、库存成本、收入确认或毛利。}
\end{exhibitbox}

钨链的政策证据比普通小金属更硬：官方口径把仲钨酸铵、氧化钨、碳化钨、固态钨、高密度钨合金以及相关生产技术资料纳入出口管制或许可框架，并说明政策目的包括国家安全、防扩散和供应链稳定。对估值的含义是可以保留资源安全溢价和供给约束敏感性，但不能机械推导“所有钨公司都能涨价”。真正能进入盈利模型的仍是公司销售价格、库存成本、产品结构、出口许可和客户订单。

政策证据之后，还需要公司层面的资源/配额兑现证据。本轮新增的上游资源安全门禁把钨矿山、钨精矿、APT、粉末加工、萤石、无水氟化氢和HFC配额放到同一张基准现金流表里。它解决的是“哪些公司有资源或配额基准可以进入估值”的问题，而不是“哪些公司已经拿到半导体单品成长订单”。因此，资源或配额可进入基准业务现金流：厦门钨业、中钨高新、章源钨业、金石资源、巨化、三美和永和可以从资源、加工或配额角度进入基准估值；翔鹭钨业更偏加工和光伏钨丝观察。所有这些证据都不替代半导体高纯钨材、WF6/NF3或含氟电子材料的客户订单、合同平均售价、收入确认比例和单品毛利。

\begin{exhibitbox}[上游资源安全与基准现金流证据]
\scriptsize
\begin{tabularx}{\exhibitboxwidth}{L{0.85cm} L{1.15cm} L{1.65cm} X X X X}
\toprule
\textbf{代码} & \textbf{名称} & \textbf{资源/配额层} & \textbf{资源或配额锚} & \textbf{产销锚} & \textbf{估值用途} & \textbf{仍缺字段} \\
\midrule
600549 & 厦门钨业 & 钨矿山+冶炼粉末+钨丝/硬质合金 & 三座生产中钨矿山和一座在建矿山合计约12,000吨钨精矿产能，钨冶炼粉末规模与品质处世界前列。 & 细钨丝销量1,292亿米，其中光伏用钨丝1,023亿米；钨粉、硬质合金和细钨丝均有产销库存表。 & 可进入钨价、资源安全和加工价差基准估值；半导体高纯钨材仍只保留外延期权。 & 半导体高纯钨材客户订单、收入占比、合同平均售价和单品毛利。 \\
000657 & 中钨高新 & 中国五矿钨产业平台 & 2025年钨精矿产量占全国8.5\%，APT占12.2\%，硬质合金产量占25.5\%；石竹园被披露为全球最大单体在产钨矿山。 & 硬质合金销售10,756吨；钨钼钽铌库存因保订单备货增加55.41\%。 & 可进入钨产业平台、资源安全和加工份额估值；高纯钨材增长每股收益仍需单品验证。 & 高纯钨材客户、订单、收入占比、合同平均售价和单品毛利。 \\
002378 & 章源钨业 & 钨矿/粉末/硬质合金一体化 & 年报披露钨粉产量在单个企业中排名第一，碳化钨粉排名第二，具备矿山、冶炼、粉末和深加工一体化基础。 & 钨粉销售4,928吨、碳化钨粉销售5,812吨；粉末端产销锚较清晰。 & 可进入粉末端竞争地位和中游加工价差估值；不能升级为半导体材料成长信用。 & 半导体高纯钨材客户订单、收入占比、合同平均售价和单品毛利。 \\
002842 & 翔鹭钨业 & 钨粉/硬质合金/光伏钨丝加工弹性 & 现有证据更偏加工和光伏钨丝，未形成可审计矿山资源自给率或半导体高纯钨材客户锚。 & 硬质合金销量同比增长70.78\%，钨丝销量同比增长217.48\%，推进300亿米超细钨丝项目。 & 可进入加工弹性和光伏钨丝观察；资源安全溢价和半导体钨材信用均需折扣。 & 资源自给率、半导体高纯钨材客户订单、收入占比、合同平均售价和单品毛利。 \\
603505 & 金石资源 & 萤石资源+无水氟化氢 & 年报披露单一萤石基本盘稳固、市场份额领先，包头选化一体项目成为全球萤石选矿和无水氟化氢领域领先产能。 & 自产萤石精矿生产41.64万吨、销售41.70万吨；自产无水氟化氢生产22.49万吨、销售21.21万吨；包钢金石萤石精粉生产80.11万吨、销售82.29万吨。 & 可进入萤石资源端和AHF成本曲线估值；液冷氟化液或电子级HF仍只给观察信用。 & 电子级HF/液冷氟化液客户、订单、收入占比、合同平均售价和单品毛利。 \\
600160 & 巨化股份 & 制冷剂配额+氟化工一体化 & 生态环境部附件命中巨化体系HFC生产配额240,826吨，全国生产配额份额30.18\%；HFC-32、HFC-134a、HFC-125和HFC-143a重点品种份额均较高。 & 制冷剂销售58.27万吨，氟化工原料利用率88.11\%、制冷剂利用率68.96\%。 & 可进入制冷剂配额现金流和氟化工龙头分部估值基准；电子湿化学品/WF6映射只能折扣。 & 电子材料收入、客户、订单、合同价格和单品毛利。 \\
603379 & 三美股份 & 制冷剂配额+AHF配套 & 生态环境部附件命中三美体系HFC生产配额76,910吨，全国生产配额份额9.64\%，并披露进口配额8,451吨二氧化碳当量。 & 氟制冷剂销售12.48万吨，R32利用率90.71\%，AHF利用率58.47\%。 & 可进入制冷剂配额和价差基准估值；半导体含氟电子材料相关性弱。 & 电子材料客户、订单、收入占比、合同平均售价和单品毛利。 \\
605020 & 永和股份 & 萤石/AHF/制冷剂一体化+HFC配额 & 生态环境部附件命中永和体系HFC生产配额60,281吨、全国生产份额7.56\%；年报披露公司为氟化工产业链较完整企业之一。 & 含氟精细化学品销量同比增长617.39\%，化工原料利用率89.46\%；HFCs配额6.03万吨位于行业前列。 & 可进入一体化氟材料和制冷剂配额基准估值；电子材料第二曲线仍需客户和订单。 & 含氟电子材料客户、订单、收入占比、合同平均售价和单品毛利。 \\
\bottomrule
\end{tabularx}
\sourcenote{2025年报产销库存、竞争格局、成本毛利、半导体钨材边界、含氟电子材料边界和生态环境部HFC配额附件派生整理；资源、配额和一体化证据只进入基准业务现金流或资源安全溢价，不替代半导体高纯钨材、WF6/NF3或含氟电子材料客户订单、合同平均售价、收入确认比例、单品收入或单品毛利。完整证据见 data/resource\_security\_evidence.json/md。}
\end{exhibitbox}

\begin{exhibitbox}[公司产品能力官方证据]
\scriptsize
\begin{tabularx}{\exhibitboxwidth}{L{1.6cm} L{1.7cm} X X X}
\toprule
\textbf{公司/层级} & \textbf{能力层} & \textbf{产品或标准细节} & \textbf{估值用途} & \textbf{证据边界} \\
\midrule
中船特气 & 公司官网产品组合 & 高纯WF6属于官网列示的主要电子特气产品；公司称具备电子特气及含氟新材料等70余种产品生产能力。 & 支持中船特气WF6/NF3直接性和调研优先级；具体估值仍受订单、平均售价、收入拆分和毛利门禁约束 & 官网产品组合不能替代客户认证、订单金额、合同价格、收入占比或单品毛利。 \\
中巨芯-U & 公司官网产品页 & 官网称WF6主要用于半导体晶圆制造成膜工艺、钨薄膜化学沉积和原子层沉积，尤其在3D结构中重要。 & 支持中巨芯具备WF6产品能力和半导体应用直接性；公告未披露长单前仍只能给观察/期权信用 & 产品页证明产品和规格，不证明产能利用、客户认证、订单、销售价格或利润。 \\
昊华科技 & 公司产品页+公告收入边界 & 有产品能力和低收入占比双重证据，适合写成能力可见但业绩弹性有限。 & 保留电子化学品平台价值，但WF6不得按高弹性主线给同等估值信用 & 产品页和0.13\%收入占比不能替代订单、平均售价、客户结构或产品毛利。 \\
和远气体 & 公司公告项目阶段 & 项目能力处在试生产/验证前后，当前不能按量产供货处理。 & 仅作为项目期权和风险提示，不进入2026E增长每股收益 & 试生产公告不能替代稳产量产、正式认证、订单金额、收入确认或毛利。 \\
电子气体WF6 & 国家标准项目 & 证明WF6进入电子气体标准化语境，可用于解释质量、纯度和产业规范化要求。 & 支持行业规范化和产品质量门槛叙事，不进入单家公司盈利模型 & 标准项目不能证明任何公司客户认证、量产订单、合同平均售价、销售收入或毛利。 \\
\bottomrule
\end{tabularx}
\sourcenote{公司官网、公司公告和全国标准信息公共服务平台归档整理；产品能力证据只证明产品组合、规格、应用或项目阶段，不替代客户认证、订单金额、合同平均售价、收入确认比例或单品毛利。}
\end{exhibitbox}

公司产品能力证据把“有无产品”和“能否兑现业绩”分开：中船特气和中巨芯-U具备更直接的WF6产品能力证据，昊华科技有产品能力但收入占比很低，和远气体仍处试生产阶段，标准化项目只说明产业规范化正在推进。能进入估值模型的仍然不是“产品存在”，而是客户认证、订单金额、合同价格、交付和毛利。

\begin{exhibitbox}[项目/产能建设官方证据]
\scriptsize
\begin{tabularx}{\exhibitboxwidth}{L{1.5cm} X X X X}
\toprule
\textbf{项目层} & \textbf{官方证据} & \textbf{产能或范围} & \textbf{估值用途} & \textbf{证据边界} \\
\midrule
投资公告 & 公司公告披露拟在邯郸市肥乡区建设年产3383吨高纯硫化氢等电子气体建设项目，预计项目总投资约86994万元，将根据项目进度分批投入。 & 新增3383吨/年电子气体（51种）生产能力，其中3158吨/年高纯电子气、225吨/年电子混合气；总周期36个月、建设周期18个月。 & 支持扩产项目存在、资本投入规模和产品多样化方向；只能提高调研优先级和项目期权权重。 & 投资公告不替代WF6单品产能拆分、产能利用率、客户认证、订单金额、合同平均售价、收入确认比例或单品毛利。 \\
环评分期与总产能 & 公司官网环评征求意见稿公示披露项目位于现有厂区内、不新增占地，项目总投资91305万元，分两期建设。 & 51种各类电子气体合计3383吨；一期工程963吨，二期工程2420吨；现有7种混合气80吨产品技术改造不新增产能。 & 支持项目分期、投资额和总产能口径，用于扩产兑现监测和资本开支解释。 & 环评总产能不替代WF6单品产能拆分、项目达产利用率、订单金额、合同平均售价、收入确认比例或单品毛利。 \\
工艺装置/车间范围 & 公司官网环评第一次公示披露建设内容包括1\#至3\#高纯气平台、八氟丙烷平台、2\#电解制氟车间、六氟化钨六氟化钼车间、尾气处理装置、钢瓶库等。 & 项目建设范围出现六氟化钨六氟化钼车间和电解制氟车间，说明项目链条覆盖电子含氟气体制造基础设施。 & 支持把项目纳入电子特气扩产和含氟工艺能力讨论，但不把车间名称直接折算成收入。 & 车间/装置范围不替代WF6单品产能拆分、试车结果、客户认证、订单金额、合同平均售价或毛利。 \\
现有工程/环保手续 & 公司官网环评第一次公示披露现有工程主要包括高纯三氟化氮、六氟化钨、三氟甲磺酸及其衍生物等特种电子气体生产装置及配套公辅工程，均已取得相关环保手续并采取污染防治措施。 & 证明现有工程已经包含高纯NF3、WF6和相关含氟电子化学品装置，且环保手续和污染防治属于已披露约束。 & 支持现有WF6/NF3工程直接性和合规基础设施讨论；仍不构成订单或价格证据。 & 现有工程披露不替代客户认证、产能利用率、订单金额、合同平均售价、收入占比或单品毛利。 \\
\bottomrule
\end{tabularx}
\sourcenote{中船特气投资公告和公司官网环评公示归档整理；项目产能证据不得替代WF6单品产能拆分、客户认证、订单金额、合同平均售价、收入确认比例或单品毛利。}
\end{exhibitbox}

项目/产能建设官方证据解决的是“扩产项目是否真实存在、投资额和装置范围是否有官方文件支持”。中船特气公告和环评公示能够证明3383吨电子气体项目、86994万元公告投资口径、91305万元环评投资口径、两期963吨/2420吨分期、六氟化钨六氟化钼车间以及现有WF6/NF3工程。它不能证明媒体口径中的WF6单品新增产能拆分，也不能证明项目已经取得客户认证、形成订单、锁定合同价格或进入收入确认。因此本报告把项目证据放入扩产监测和事件期权，不把它直接折成2026E增长每股收益。

\section{核心覆盖池}
覆盖池按产业链位置而非概念板块划分。直接WF6/NF3产能公司的估值弹性最高，但证伪速度也最快；钨矿和制冷剂公司的利润确认更直接；下游晶圆厂和设备厂只作为需求验证，不在本报告内给同等深度目标价。

\begin{exhibitbox}[18只核心标的产业链映射]
\scriptsize
\resizebox{\exhibitwidth}{!}{
\begin{tabular}{lllll}
\toprule
\textbf{代码} & \textbf{名称} & \textbf{环节} & \textbf{角色} & \textbf{证据} \\
\midrule
000657 & 中钨高新 & 硬质合金/数控刀具 & 钨下游硬质合金和刀具平台，价格弹性低于矿端但制造升级属性强 & 中 \\
002378 & 章源钨业 & 钨矿/冶炼/深加工 & 钨产业链完整公司，资源储量与钨价弹性较高 & 高 \\
002407 & 多氟多 & 氟化盐/电子级HF/锂电材料 & 电子级氢氟酸与氟化盐平台，受锂电周期扰动 & 中 \\
002549 & 凯美特气 & 稀有气体/工业尾气回收 & 稀有气体和尾气回收，受氦氖氪氙及电子气体情绪影响 & 中 \\
002842 & 翔鹭钨业 & APT/钨粉/硬质合金 & 钨加工弹性标的，受益补库但利润受原料波动影响 & 中 \\
002971 & 和远气体 & 电子级氯气/氯化氢/园区供气 & 潜江电子特气园区与电子级氯气、氯化氢等品类扩张 & 中 \\
300037 & 新宙邦 & 有机氟/电子氟化液/电解液 & 有机氟精细化学品和电子氟化液，AI数据中心液冷为远期映射 & 中 \\
300346 & 南大光电 & 电子特气/前驱体/光刻胶 & 磷烷、砷烷等特气与前驱体材料平台，WF6直接性弱于中船/中巨芯 & 中 \\
600160 & 巨化股份 & 制冷剂/氟化工/中巨芯股权 & 制冷剂配额龙头，并通过中巨芯映射电子级HF和含氟电子特气 & 高 \\
600378 & 昊华科技 & 电子特气/含氟材料/科研平台 & 具备WF6产品和含氟材料平台，但公司公告显示WF6收入占比低 & 高 \\
600549 & 厦门钨业 & 钨资源/冶炼/硬质合金 & 钨一体化与高端钨材龙头，受益钨价与战略资源重估 & 高 \\
603379 & 三美股份 & 制冷剂 & 三代制冷剂盈利弹性强，半导体属性弱于巨化/昊华 & 高 \\
603505 & 金石资源 & 萤石/氟化工原料 & 萤石上游，间接受益制冷剂和含氟材料景气 & 中 \\
605020 & 永和股份 & 含氟材料/制冷剂 & 氟化工产业链一体化，受益制冷剂与含氟材料景气 & 中 \\
688106 & 金宏气体 & 大宗+电子气体 & 综合气体服务，电子特气成长但短期利润基数偏弱 & 中 \\
688146 & 中船特气 & WF6/NF3电子特气 & 国内电子特气龙头，公告披露WF6现有产能2000吨/年、6N级 & 高 \\
688268 & 华特气体 & 电子特气平台 & 电子特气品类平台和光刻气国产替代，WF6为品类扩张映射 & 高 \\
688549 & 中巨芯-U & 电子级HF/WF6/含氟电子特气 & 高纯WF6产能600吨但公告提示无新长期/大额订单、无扩产计划 & 高 \\
\bottomrule
\end{tabular}
}
\end{exhibitbox}

上表只回答“公司在链条哪里”，还不能回答“能否进入估值”。真正决定估值信用的是关系证据矩阵：收入暴露是否披露、产能和认证是否可核验、订单是否可见、缺口是否会影响利润确认。电子特气标的中，中船特气的WF6产能和6N规格直接性最高，且已有WF6/NF3历史合同金额样本；但仍缺合同平均售价、收入确认比例、客户份额、单品收入和单品毛利，以及2026年以后持续订单。中巨芯-U具备高纯WF6产能，但公告边界显示暂无新增长期或大额订单；昊华科技有产品能力，但WF6收入占比低；华特气体则明确无WF6合成产能，仅按需供应纯化产品且占比较小。扩散型气体公司若不能给出产品级收入或客户证明，只能保留观察席位。

\
                \begin{exhibitbox}[钨链公司关系证据]
                \footnotesize
\begin{tabularx}{\exhibitboxwidth}{L{0.9cm} L{1.2cm} L{1.3cm} X X X}
\toprule
\textbf{代码} & \textbf{名称} & \textbf{直接性} & \textbf{收入/产品暴露} & \textbf{主要缺口} & \textbf{估值使用} \\
\midrule
002842 & 翔鹭钨业 & 直接 & 年报产销/产能证据：硬质合金产/销+84.67\%/+70.78\%；钨丝产/销+75.06\%/+217.48\%。；竞争地位：加工和光伏钨丝弹性强，300亿米超细钨丝项目提供成长观察点。 & 高纯钨材客户、订单、收入占比和单品毛利仍缺 & 资源/周期证据可用于估值 \\
002378 & 章源钨业 & 直接 & 年报产销/产能证据：钨粉销4,928吨；碳化钨粉销5,812吨；硬质合金库存下降。；竞争地位：粉末端地位强：钨粉单企第一、碳化钨粉第二。 & 高纯钨材客户、订单、收入占比和单品毛利仍缺 & 资源/周期证据可用于估值 \\
600549 & 厦门钨业 & 直接 & 年报产销/产能证据：细钨丝销1,292亿米；钨粉/硬质合金产销库存完整。；竞争地位：全链钨龙头：矿山约12,000吨产能，钨丝份额前列，硬质合金规模国内前列。 & 高纯钨材客户、订单、收入占比和单品毛利仍缺 & 资源/周期证据可用于估值 \\
000657 & 中钨高新 & 中等 & 年报产销/产能证据：硬质合金销10,756吨；钨钼钽铌库存因保订单备货增55.41\%。；竞争地位：钨平台龙头：钨精矿8.5\%、APT12.2\%、硬质合金25.5\%，全产业链强。 & 高纯钨材客户、订单、收入占比和单品毛利仍缺 & 资源/周期证据可用于估值 \\
\bottomrule
\end{tabularx}
\end{exhibitbox}

\
                \begin{exhibitbox}[氟化工公司关系证据]
                \footnotesize
\begin{tabularx}{\exhibitboxwidth}{L{0.9cm} L{1.2cm} L{1.3cm} X X X}
\toprule
\textbf{代码} & \textbf{名称} & \textbf{直接性} & \textbf{收入/产品暴露} & \textbf{主要缺口} & \textbf{估值使用} \\
\midrule
603379 & 三美股份 & 直接 & 制冷剂/氟化工现金流；电子材料收入仍需披露 & 缺公司级电子材料客户、订单和收入占比 & 盈利支撑与期权信用拆分后可用于估值 \\
300037 & 新宙邦 & 中等 & 制冷剂/氟化工现金流；电子材料收入仍需披露 & 缺公司级电子材料客户、订单和收入占比 & 盈利支撑与期权信用拆分后可用于估值 \\
605020 & 永和股份 & 直接 & 制冷剂/氟化工现金流；电子材料收入仍需披露 & 缺公司级电子材料客户、订单和收入占比 & 盈利支撑与期权信用拆分后可用于估值 \\
600160 & 巨化股份 & 直接+映射 & 制冷剂/氟化工现金流；电子材料收入仍需披露 & 缺公司级电子材料客户、订单和收入占比 & 盈利支撑与期权信用拆分后可用于估值 \\
002407 & 多氟多 & 中等 & 制冷剂/氟化工现金流；电子材料收入仍需披露 & 缺公司级电子材料客户、订单和收入占比 & 盈利支撑与期权信用拆分后可用于估值 \\
603505 & 金石资源 & 间接 & 制冷剂/氟化工现金流；电子材料收入仍需披露 & 缺公司级电子材料客户、订单和收入占比 & 盈利支撑与期权信用拆分后可用于估值 \\
\bottomrule
\end{tabularx}
\end{exhibitbox}

\
                \begin{exhibitbox}[电子特气公司关系证据]
                \footnotesize
\begin{tabularx}{\exhibitboxwidth}{L{0.9cm} L{1.2cm} L{1.3cm} X X X}
\toprule
\textbf{代码} & \textbf{名称} & \textbf{直接性} & \textbf{收入/产品暴露} & \textbf{主要缺口} & \textbf{估值使用} \\
\midrule
600378 & 昊华科技 & 产品直接/收入低 & 有WF6产品能力；2025年WF6收入占比0.13\% & 收入占比低，缺订单、价格和客户结构 & 仅事件驱动或观察信用 \\
300346 & 南大光电 & 平台映射 & 电子特气平台或扩散暴露；WF6直接收入未找到 & 缺具名客户、订单、价格和产品收入占比 & 仅事件驱动或观察信用 \\
688146 & 中船特气 & 直接 & WF6/NF3产能和规格直接；部分客户洽谈增量订单 & 缺具名客户、订单金额、平均售价和产品毛利 & 仅事件驱动或观察信用 \\
688549 & 中巨芯-U & 直接 & 高纯WF6产能直接；盈利基数和收入结构未披露 & 无新增长期/大额订单，缺客户和价格 & 仅事件驱动或观察信用 \\
688268 & 华特气体 & 平台映射 & 电子特气平台；无WF6合成产能，仅纯化供应且占比较小 & 缺客户名称、订单金额、收入占比和产品毛利 & 仅事件驱动或观察信用 \\
002971 & 和远气体 & 品类扩张 & 电子特气平台或扩散暴露；WF6直接收入未找到 & 缺具名客户、订单、价格和产品收入占比 & 仅事件驱动或观察信用 \\
002549 & 凯美特气 & 扩散映射 & 电子特气平台或扩散暴露；WF6直接收入未找到 & 缺具名客户、订单、价格和产品收入占比 & 仅事件驱动或观察信用 \\
688106 & 金宏气体 & 扩散映射 & 电子特气平台或扩散暴露；WF6直接收入未找到 & 缺具名客户、订单、价格和产品收入占比 & 仅事件驱动或观察信用 \\
\bottomrule
\end{tabularx}
\end{exhibitbox}


氟化工链条必须先拆成“制冷剂现金流”和“含氟电子材料期权”。生态环境部配额可以证明巨化、三美、永和体系的制冷剂供给约束和相对份额；主营构成、产销库存和成本毛利证据可以解释氟化工主业韧性；但这些证据都不能直接证明电子级HF、氟化液、有机氟电子材料或含氟电子特气已经进入具名半导体客户、形成订单、锁定价差并贡献单品毛利。

\begin{exhibitbox}[含氟电子材料边界]
\scriptsize
\begin{tabularx}{\exhibitboxwidth}{L{0.85cm} L{1.15cm} X X X X X}
\toprule
\textbf{代码} & \textbf{名称} & \textbf{基准现金流锚} & \textbf{电子材料相关性} & \textbf{客户/订单} & \textbf{估值用途} & \textbf{仍缺字段} \\
\midrule
600160 & 巨化股份 & 生态环境部附件命中巨化体系HFC生产配额240,826吨、全国生产份额30.18\%；年报/主营构成支持制冷剂和氟化工主业现金流。 & 通过中巨芯股权映射电子级HF和含氟电子特气，但上市公司口径仍以制冷剂/氟化工现金流为主。 & 未披露巨化本体含氟电子材料具名客户、批量订单金额、客户采购量或收入确认比例。 & 制冷剂配额和价差进入基准现金流与分部估值主业权重；中巨芯电子材料映射只给股权/期权信用。 & 含氟电子材料客户、订单金额、单品价差、收入占比和单品毛利。 \\
603379 & 三美股份 & 生态环境部附件命中三美体系HFC生产配额76,910吨、全国生产份额9.64\%，并披露进口配额8,451吨二氧化碳当量。 & 当前公司逻辑更纯粹指向制冷剂配额和价差，半导体含氟电子材料相关性弱于巨化/中巨芯链条。 & 未披露电子级HF、氟化液或含氟电子材料客户认证、订单金额或交付节奏。 & 以制冷剂配额现金流和价差弹性估值；不得套用电子材料高成长分部倍数。 & 含氟电子材料客户、订单、价差、收入占比和单品毛利。 \\
605020 & 永和股份 & 生态环境部附件命中永和体系HFC生产配额60,281吨、全国生产份额7.56\%；年报产销证据支持一体化氟化工主业。 & 年报和主营证据支持含氟材料/制冷剂一体化平台，但未把含氟电子材料单品与半导体客户或订单绑定。 & 未披露含氟电子材料具名客户、认证阶段、订单金额、客户份额或收入确认比例。 & 一体化氟化工和制冷剂配额进入基准业务；电子材料第二曲线只保留调研信用。 & 含氟电子材料客户、订单、单品价差、收入占比和单品毛利。 \\
002407 & 多氟多 & 生态环境部HFC附件未检出多氟多同名配额行；主营锚来自氟化盐、电子级HF、锂电材料和新材料平台。 & 电子级HF和氟基新材料具备主题相关性，但锂电材料周期和电子化学品认证需要拆开处理。 & 未披露电子级HF/含氟电子材料具名半导体客户、订单金额、单品交付节奏或确认比例。 & 按氟化盐、电子化学品和锂电材料平台做分部估值折扣；不得按制冷剂配额龙头或电子材料高成长段直接上修。 & 电子级HF/含氟电子材料客户、订单金额、单品价差、收入占比和单品毛利。 \\
300037 & 新宙邦 & 主营锚来自电解液、有机氟精细化学品和电子氟化液平台，公告扫描命中锂电电解液合作协议而非含氟电子材料订单。 & 有机氟和电子氟化液可映射AI数据中心液冷和半导体材料，但当前公开证据更多是远期应用线索。 & 未披露电子氟化液或含氟电子材料具名客户、批量供货订单、单品客户份额或收入确认比例。 & 锂电电解液和有机氟平台分部可进入基准/分部估值；电子氟化液第二曲线只给观察信用。 & 电子氟化液/含氟电子材料客户、订单、单品价差、收入占比和单品毛利。 \\
603505 & 金石资源 & 年报和主营锚支持萤石资源、无水氟化氢和上游氟化工原料暴露。 & 处于萤石/HF上游，受制冷剂和含氟材料景气传导；不是电子级HF或含氟电子材料订单直接主体。 & 未披露含氟电子材料客户、电子级HF订单、客户份额或收入确认比例。 & 按萤石资源和氟化工原料景气估值；不能使用电子材料第二曲线增长倍数。 & 电子级HF/含氟电子材料客户、订单、收入占比、单品价差和毛利。 \\
\bottomrule
\end{tabularx}
\sourcenote{生态环境部HFC配额附件、主营构成、产销量/库存、成本毛利、竞争格局、合同公告和年报主题证据归档整理；该表只证明制冷剂/氟化工基准现金流与电子材料第二曲线的边界，不替代含氟电子材料客户、订单、单品价差、收入占比或单品毛利。完整证据见 data/fluorinated\_electronic\_material\_boundary.json/md。}
\end{exhibitbox}

当前边界是0/6可进入含氟电子材料增长每股收益。巨化、三美、永和可以把制冷剂配额、价差和一体化氟化工现金流放入基准估值；多氟多、新宙邦、金石资源可以分别保留电子级HF/新材料、有机氟/电子氟化液、萤石/HF上游的观察价值。它们共同的硬约束是：未披露具名客户、订单金额、单品价差、收入占比和单品毛利前，含氟电子材料第二曲线不能进入基本面目标价上修。


钨链年报证据已经能把“钨价弹性”拆成公司级兑现能力，而不是只停留在题材标签。厦门钨业披露细钨丝销量1,292亿米，其中光伏用钨丝1,023亿米，并有钨粉、硬质合金和细钨丝产销库存数据；中钨高新披露硬质合金销售10,756吨、钨钼钽铌库存因保订单备货增加55.41\%，同时具备钨精矿、APT和硬质合金产量份额证据；章源钨业披露钨粉销售4,928吨、碳化钨粉销售5,812吨，且钨粉单企产量第一、碳化钨粉第二；翔鹭钨业披露硬质合金销量同比增长70.78\%、钨丝销量同比增长217.48\%，并推进300亿米超细钨丝项目。这些证据可以支持资源/加工价差、光伏钨丝和硬质合金制造升级的基本面讨论，但仍不得替代半导体高纯钨材客户、订单金额、合同平均售价、收入占比或单品毛利。

\begin{exhibitbox}[半导体高纯钨材边界]
\scriptsize
\begin{tabularx}{\exhibitboxwidth}{L{0.85cm} L{1.15cm} X X X X X}
\toprule
\textbf{代码} & \textbf{名称} & \textbf{年报/主营锚} & \textbf{半导体相关性} & \textbf{客户/订单} & \textbf{估值用途} & \textbf{仍缺字段} \\
\midrule
600549 & 厦门钨业 & 钨钼等有色金属产品收入占比42.4\%、毛利率28.4\%；钨粉末销量9,670吨，硬质合金销量8,272吨，细钨丝销量1,292亿米，其中光伏用钨丝1,023亿米。 & 年报语境提及半导体领域六氟化钨增量贡献和高端靶材应用，但公司层面披露仍主要落在钨粉、硬质合金、细钨丝和钨钼产品大类。 & 未披露半导体高纯钨材客户、认证阶段、订单金额或收入确认比例。 & 钨资源/加工主业可进入周期和资源安全估值；半导体钨材只能作为外延期权和调研优先级。 & 半导体高纯钨材客户、订单金额、收入占比、合同平均售价和单品毛利。 \\
000657 & 中钨高新 & 硬质合金及钨产品收入占比100.0\%、毛利率23.7\%；硬质合金销售10,756吨，钨钼钽铌库存因保证订单执行增加55.41\%。 & 年报披露钨粉可用于钨丝、钨棒、钨管等钨材及钨合金，应用场景包括集成电路和半导体等工业领域；互动问答线索显示没有进入WF6计划，但存在向相关客户供应钨粉原料的链条线索。 & 未披露半导体高纯钨材具名客户、认证阶段、订单金额、订单期限或确认比例。 & 制造升级和刀具平台可进入基准估值；半导体高纯钨材不得给成长每股收益信用。 & 高纯钨材客户/订单、产品收入占比、合同平均售价和产品级毛利。 \\
002378 & 章源钨业 & 碳化钨粉收入占比35.7\%、毛利率11.5\%；钨粉收入占比31.0\%、毛利率8.5\%；钨粉销量4,928吨，碳化钨粉销量5,812吨。 & 年报主题扫描未找到半导体关键词；披露ZWC02超细碳化钨粉和超粗碳化钨粉通过客户认证，但没有把认证映射到半导体高纯钨材客户或订单。 & 未披露半导体高纯钨材客户、订单金额、客户份额、供货节奏或收入确认比例。 & 资源冶炼一体化和钨粉景气可进入周期/加工估值；半导体材料外延仍作折扣期权。 & 半导体高纯钨材客户/订单、单品收入、合同平均售价和单品毛利。 \\
002842 & 翔鹭钨业 & 粉末制品收入占比66.9\%、毛利率12.9\%；硬质合金收入占比20.8\%、毛利率15.5\%；钨丝系列收入占比8.2\%、毛利率24.0\%。 & 年报把半导体作为宏观需求场景提及；公司推进年产300亿米超细钨丝项目，当前证据更多对应光伏金刚线母线和硬质合金放量。 & 互动问答显示光伏金刚线用钨丝订单和意向订单充足，但不是半导体高纯钨材订单；未披露半导体客户认证或订单金额。 & 钨加工、硬质合金和光伏钨丝兑现可进入基准估值；不得把光伏钨丝项目外推为半导体钨材成长每股收益。 & 半导体高纯钨材客户、订单、产品收入、合同平均售价和单品毛利。 \\
\bottomrule
\end{tabularx}
\sourcenote{2025年年报、主营构成、产销量/库存、投资者关系互动问答和年报短摘归档整理；该表只证明钨链资源/加工和半导体外延的边界，不替代半导体高纯钨材客户、订单、合同平均售价、收入占比或单品毛利。完整证据见 data/semiconductor\_tungsten\_material\_boundary.json/md。}
\end{exhibitbox}

钨链四家公司已经能解释资源/加工利润，但还不能解释半导体高纯钨材成长利润。当前边界是0/4可进入半导体高纯钨材增长每股收益：厦门钨业的钨粉、硬质合金和细钨丝产销完整，中钨高新的硬质合金和刀具平台较强，章源钨业有钨粉和碳化钨粉大类收入，翔鹭钨业有粉末制品、硬质合金和光伏钨丝放量；这些证据可以进入钨价、资源安全、库存和加工价差讨论。它们不能替代半导体客户认证、订单金额、合同平均售价、收入占比和单品毛利，因此半导体高纯钨材外延仍只保留期权信用。



官方年报客户链已经把“具名客户未找到”的问题收窄为“单品份额和订单未披露”。中船特气年报披露NF3稳定供应台积电、美光、海力士、英飞凌、格罗方德、中芯国际、长江存储、华虹、京东方等客户，WF6稳定供应台积电、美光、海力士、英飞凌、铠侠、格罗方德、中芯国际、长江存储、长鑫存储、华虹等客户；中巨芯-U披露电子级氢氟酸、硫酸、硝酸等已在中芯国际、华虹集团等客户批量供货，电子特气及前驱体进入中芯国际、华虹集团、SK海力士、日本铠侠、上海华力、士兰微等试用与供应阶段；华特气体披露获得中芯国际、长江存储、三星、SK海力士、台积电等客户认可；金宏气体披露服务中芯国际、海力士、长鑫存储等头部企业；南大光电披露前驱体导入国内领先芯片制造企业量产制程。这些证据能提升客户链置信度和调研优先级，但仍不得替代WF6/NF3、电子级HF、前驱体等单品客户份额、订单金额、合同平均售价、收入确认比例或单品毛利。


客户侧验证探针进一步把“供应商说供给了谁”和“客户侧是否能反向验证采购”分开。本轮8个锚点中，供应商侧年报、认证或平台需求证据较强，但客户侧订单、采购清单、供应商名录、采购份额或合同价格文件均未取得，客户侧反向验证为0/8。因此，客户链证据可以提高关系置信度和下一轮调研优先级，不能直接变成订单金额、增长收入或目标价上修。

\begin{exhibitbox}[客户侧验证探针]
\scriptsize
\begin{tabularx}{\exhibitboxwidth}{L{1.65cm} L{1.45cm} L{1.75cm} X X X}
\toprule
\textbf{客户链锚点} & \textbf{关联标的} & \textbf{状态} & \textbf{客户侧证据} & \textbf{估值用途} & \textbf{仍缺字段} \\
\midrule
中船特气WF6/NF3客户链 & 688146 中船特气 & 供应商侧具名客户强；客户侧采购文件未取得 & 本报告语料内未取得上述晶圆厂或存储厂客户侧供应商名录、采购订单、采购份额或合同价格文件。 & 关系置信度和调研优先级可提高；客户侧文件未取得前，不把客户名单直接折成WF6/NF3订单、收入或增长每股收益。 & 客户侧供应商名录、采购订单、采购份额、合同平均售价、收入确认比例和单品毛利。 \\
中巨芯电子湿化学品/电子特气客户链 & 688549 中巨芯-U；600160 巨化股份 & 供应商侧批量供货/试用披露强；客户侧订单金额未取得 & 本报告语料内未取得中芯国际、华虹、SK海力士、铠侠、上海华力或士兰微客户侧采购清单、产品份额或订单金额文件。 & 可进入强观察池和客户链置信度；WF6高成长信用仍受无新增长期/大额订单和无客户侧采购文件约束。 & WF6客户侧订单、电子级HF客户采购量、产品收入占比、合同平均售价和单品毛利。 \\
华特气体认证客户链 & 688268 华特气体 & 认证/客户覆盖证据强；WF6纯化客户侧订单未取得 & 本报告语料内未取得客户侧采购订单、WF6纯化客户名单、纯化产品份额或客户采购价格文件。 & 支持平台质量和认证壁垒；无WF6合成产能且纯化业务占比较小时，不给WF6单品增长每股收益。 & 纯化WF6客户侧采购、收入占比、订单金额、合同平均售价和产品毛利。 \\
南大光电前驱体客户链 & 300346 南大光电 & 供应商侧量产导入披露强；客户侧反向文件未取得 & 本报告语料内未取得国内领先芯片制造企业客户侧供应商名录、采购份额、价格或订单金额文件。 & 可支持强观察池和材料平台溢价；客户侧文件未取得前，不把前驱体导入直接折成订单金额或单品毛利。 & 具名客户、采购份额、订单金额、合同平均售价、收入确认比例和单品毛利。 \\
金宏气体半导体客户链 & 688106 金宏气体 & 客户导入证据存在；客户侧采购细节未取得 & 本报告语料内未取得中芯国际、海力士、长鑫存储客户侧气体供应商名录、具体电子气体采购量或订单价格文件。 & 支持电子气体平台观察信用；工业现场制气合同不能迁移为电子特气单品订单。 & 电子特气单品客户、采购量、订单金额、合同平均售价和产品毛利。 \\
凯美特气光刻气认证链 & 002549 凯美特气 & 认证证据强于订单证据；客户侧采购文件未取得 & 本报告语料内未取得Coherent、Cymer、GIGAPHOTON或终端晶圆厂客户侧采购订单、价格或份额文件。 & 支持认证壁垒和期权信用；客户侧订单与价格未取得前，不进入增长每股收益。 & 光刻气客户侧采购订单、采购份额、收入占比、合同平均售价和产品毛利。 \\
和远气体半导体/面板认证链 & 002971 和远气体 & 认证启动和负面门禁并存；客户侧确认未取得 & 本报告语料内未取得半导体或面板客户侧认证通过文件、采购订单、销售价格或收入确认材料。 & 只能作为事件验证池；WF6不得进入2026E增长每股收益。 & 客户侧认证通过、正式订单、收入确认比例、合同平均售价和单品毛利。 \\
AI/HBM/云端算力平台链 & 全体电子特气、钨链和氟化工观察池 & 平台需求强；A股上游客户侧分配证据未取得 & 本报告语料内未取得平台厂商或终端晶圆厂把上述需求明确分配到A股上游供应商的采购订单、供应商名录或价格文件。 & 只提高调研优先级和需求锚权重；不得直接上修A股供应商基本面目标价。 & 客户侧供应商名录、订单金额、采购量、合同平均售价、产品收入和产品毛利。 \\
\bottomrule
\end{tabularx}
\sourcenote{客户侧验证探针由客户认证/具名客户披露、客户/供应商明细、合同公告、公开券商全文、下游需求和AI平台锚派生；当前客户侧订单/采购/份额反向验证为0/8。供应商年报自述、平台需求或券商关键词不替代客户侧采购文件、订单金额、合同平均售价、收入确认比例或单品毛利。完整证据见 data/customer\_side\_verification\_probe.json/md。}
\end{exhibitbox}

这个探针会直接约束估值语言：中船特气、中巨芯-U、华特气体、南大光电、金宏气体、凯美特气和和远气体可以按“供应商侧披露较强、客户侧反向文件未取得”处理；AI/HBM和云端算力平台链只能说明终端需求强，不能说明A股上游供应商已经拿到客户侧采购份额。后续若取得客户侧供应商名录、采购订单、采购份额、合同平均售价或收入确认材料，再把对应标的从观察信用升级为订单信用。


下游客户公开文件探针把反向验证再推进一层：本轮实际检查中芯国际、华虹、士兰微、中微公司、北方华创、沪硅产业、京东方A和TCL科技等客户侧2025年公开年报，扫描A股上游供应商名、WF6/NF3/电子特气/电子化学品等产品词以及采购、订单、供应商名录、价格等硬字段。结果是，客户侧公开年报已检查8/8份，但产品级采购订单、采购份额或合同价格反向验证仍为0/8。这意味着供应商侧客户名和客户侧年报产品词仍只能作为调研入口，不能写成订单和利润。

\begin{exhibitbox}[下游客户公开文件探针]
\scriptsize
\begin{tabularx}{\exhibitboxwidth}{L{1.35cm} L{1.35cm} L{1.65cm} X X X}
\toprule
\textbf{客户侧主体} & \textbf{角色} & \textbf{已检查文件} & \textbf{供应商/产品命中} & \textbf{反向验证} & \textbf{估值用途} \\
\midrule
688981 中芯国际 & 晶圆制造客户侧 & 公开年报已检查；2026-03-27 & 供应商名候选命中：巨化股份；产品词未命中；采购/订单词采购订单、合格供应商、供应商、采购 & 未形成客户侧采购订单/采购份额/合同价格反向验证 & 只用于客户侧穷尽和负面门禁；未取得产品级采购订单、份额、价格前，不进入增长每股收益或基本面目标价上修。 \\
688347 华虹公司 & 晶圆制造客户侧 & 公开年报已检查；2026-05-15 & 未命中A股上游供应商名；产品词未命中；采购/订单词采购订单、合格供应商、供应商、采购 & 未形成客户侧采购订单/采购份额/合同价格反向验证 & 只用于客户侧穷尽和负面门禁；未取得产品级采购订单、份额、价格前，不进入增长每股收益或基本面目标价上修。 \\
600460 士兰微 & IDM/功率半导体客户侧 & 公开年报已检查；2026-04-25 & 未命中A股上游供应商名；产品词未命中；采购/订单词合同价格、供应商、采购、订单 & 未形成客户侧采购订单/采购份额/合同价格反向验证 & 只用于客户侧穷尽和负面门禁；未取得产品级采购订单、份额、价格前，不进入增长每股收益或基本面目标价上修。 \\
688012 中微公司 & 半导体设备客户侧 & 公开年报已检查；2026-03-31 & 供应商名候选命中：南大光电；产品词未命中；采购/订单词供应商名录、合格供应商、供应商、采购 & 未形成客户侧采购订单/采购份额/合同价格反向验证 & 只用于客户侧穷尽和负面门禁；未取得产品级采购订单、份额、价格前，不进入增长每股收益或基本面目标价上修。 \\
002371 北方华创 & 半导体设备客户侧 & 公开年报已检查；2026-04-18 & 未命中A股上游供应商名；产品词前驱体；采购/订单词采购金额、供应商、采购、订单 & 未形成客户侧采购订单/采购份额/合同价格反向验证 & 只用于客户侧穷尽和负面门禁；未取得产品级采购订单、份额、价格前，不进入增长每股收益或基本面目标价上修。 \\
688126 沪硅产业 & 硅片材料客户侧 & 公开年报已检查；2026-04-17 & 供应商名候选命中：中巨芯、南大光电；产品词前驱体；采购/订单词供应商名录、供应商名单、合格供应商、供应商 & 未形成客户侧采购订单/采购份额/合同价格反向验证 & 只用于客户侧穷尽和负面门禁；未取得产品级采购订单、份额、价格前，不进入增长每股收益或基本面目标价上修。 \\
000725 京东方A & 面板客户侧 & 公开年报已检查；2026-04-01 & 未命中A股上游供应商名；产品词未命中；采购/订单词采购金额、合同价格、供应商、采购 & 未形成客户侧采购订单/采购份额/合同价格反向验证 & 只用于客户侧穷尽和负面门禁；未取得产品级采购订单、份额、价格前，不进入增长每股收益或基本面目标价上修。 \\
000100 TCL科技 & 面板客户侧 & 公开年报已检查；2026-03-28 & 未命中A股上游供应商名；产品词未命中；采购/订单词采购金额、合同价格、供应商、采购 & 未形成客户侧采购订单/采购份额/合同价格反向验证 & 只用于客户侧穷尽和负面门禁；未取得产品级采购订单、份额、价格前，不进入增长每股收益或基本面目标价上修。 \\
\bottomrule
\end{tabularx}
\sourcenote{巨潮资讯下游客户2025年年度报告PDF归档与pdftotext抽取；该探针检查供应商名、产品词和采购词候选，当前产品级采购订单/采购份额/合同价格反向验证为0。完整证据见 data/downstream\_customer\_public\_file\_probe.json/md。}
\end{exhibitbox}

\begin{exhibitbox}[下游需求锚行情与财务]
\scriptsize
\begin{tabularx}{\exhibitboxwidth}{L{0.9cm} L{1.2cm} L{1.4cm} L{1.5cm} L{1.2cm} L{1.4cm} X}
\toprule
\textbf{代码} & \textbf{名称} & \textbf{环节} & \textbf{一季度收入/同比} & \textbf{成交额} & \textbf{需求信号} & \textbf{材料链意义} \\
\midrule
688981 & 中芯国际 & 晶圆制造 & 176.2亿元/8.1\% & 161.7亿元 & 收入低速增长；盈利为正；成交活跃 & 晶圆代工扩产与先进/成熟制程投片，验证WF6/NF3/刻蚀清洗气体需求 \\
688347 & 华虹公司 & 晶圆制造 & 46.3亿元/18.2\% & 72.0亿元 & 收入稳健增长；盈利为正；成交活跃 & 成熟制程、功率和嵌入式非易失存储投片，验证电子特气国产替代需求 \\
688012 & 中微公司 & 半导体设备 & 29.1亿元/34.1\% & 131.5亿元 & 收入高增长；盈利为正；成交活跃 & 刻蚀/薄膜设备订单与装机，验证高纯特气在设备侧的工艺位置 \\
002371 & 北方华创 & 半导体设备 & 103.2亿元/25.8\% & 100.1亿元 & 收入高增长；盈利为正；成交活跃 & PVD/CVD/刻蚀设备国产化和扩产，验证前道材料放量条件 \\
688126 & 沪硅产业 & 硅片材料 & 10.8亿元/35.2\% & 92.2亿元 & 收入高增长；盈利亏损；成交活跃 & 硅片材料景气和晶圆厂扩产节奏，辅助判断前道材料需求周期 \\
\bottomrule
\end{tabularx}
\end{exhibitbox}

下游需求锚的意义是确认半导体前道材料国产替代的方向，而不是确认任何单一上游公司的收入。中芯国际、华虹、中微公司、北方华创和沪硅产业的收入增速、成交活跃度和利润状态，可以帮助判断晶圆制造、设备装机和硅片景气是否支撑CVD、刻蚀、清洗和薄膜材料需求；它们不能替代中船特气、中巨芯-U或昊华科技自己的客户、订单和价格披露。这个边界如果不写清楚，产业链章节就会退化成概念股票表。

\begin{exhibitbox}[AI平台/HBM/网络需求锚]
\scriptsize
\begin{tabularx}{\exhibitboxwidth}{L{1.2cm} L{1.8cm} L{1.7cm} X X X}
\toprule
\textbf{厂商} & \textbf{平台} & \textbf{平台组} & \textbf{材料链传导} & \textbf{语料缺口} & \textbf{估值用途} \\
\midrule
NVIDIA & Blackwell / Blackwell Ultra & AI GPU/HBM & 先进逻辑晶圆、CoWoS/先进封装、HBM、CVD/刻蚀/清洗特气和前道薄膜材料需求锚。 & 未取得NVIDIA客户侧材料供应商名录或A股供应商采购订单。 & 需求锚和市场叙事强度；不进入单公司每股收益。 \\
SK hynix & HBM3E / HBM4 & HBM与AI存储 & DRAM/HBM晶圆、TSV/堆叠封装、清洗刻蚀气体、前驱体和高纯材料需求锚。 & 未取得SK hynix客户侧电子特气/钨材采购清单。 & 解释HBM/3D存储需求锚；不证明国内供应商订单。 \\
Google Cloud & Cloud TPU v6e / Trillium & 云端ASIC & 定制ASIC先进制程、HBM/高速存储、先进封装和前道材料需求锚。 & 未取得Google TPU供应链中A股材料供应商认证或采购文件。 & 作为云ASIC需求锚；不进入上游材料公司收入桥。 \\
AWS & Trainium / Trainium2 & 云端ASIC与HBM & 云厂商自研ASIC、HBM、先进封装、前道制造和网络互连材料需求锚。 & 未取得AWS Trainium供应链中A股电子材料订单。 & 作为云ASIC和HBM需求锚；不进入单公司每股收益。 \\
Intel & Gaudi 3 & AI加速器与HBM & AI加速器晶圆制造、HBM、封装和高速互连材料需求锚。 & 未取得Intel Gaudi供应链映射到A股电子材料供应商的订单证据。 & 作为非NVIDIA AI加速器需求锚；不进入A股材料公司目标价。 \\
Huawei & Ascend / Atlas 900 A3 SuperPoD & 国产AI算力 & 国产AI芯片、国产晶圆/封装、国产材料替代和供应链安全需求锚。 & 未取得华为/昇腾客户侧材料采购、电子特气供应商或A股材料订单文件。 & 国产算力需求锚；只提升观察优先级，不给订单信用。 \\
NVIDIA & Spectrum-X Ethernet / Photonics & AI networking/optical & 高速交换芯片、硅光/光模块、先进封装和数据中心互连需求锚，间接影响前道材料和设备需求。 & 未取得网络/光模块需求向WF6或A股电子特气订单的直接映射。 & 作为AI集群网络需求锚；不能替代半导体材料订单。 \\
\bottomrule
\end{tabularx}
\sourcenote{NVIDIA、SK hynix、Google Cloud、AWS、Intel、Huawei公开资料整理；平台链只解释AI/HBM/云ASIC/国产算力/网络需求锚，不证明A股上游客户、订单、平均售价、收入确认或单品毛利。}
\end{exhibitbox}

平台链把“AI叙事”拆成三种可跟踪需求：第一是NVIDIA Blackwell、Google TPU、AWS Trainium和Intel Gaudi代表的AI加速器投片与HBM需求；第二是SK hynix HBM3E/HBM4代表的存储堆叠和先进封装需求；第三是华为昇腾和Spectrum-X代表的国产算力与AI网络扩容。它能解释为什么前道材料被重新定价，但不能替代供应商认证、采购订单和产品毛利。对A股标的而言，平台链只提高调研优先级，不能直接上调目标价。

\begin{exhibitbox}[公司澄清/风险提示证据]
\scriptsize
\begin{tabularx}{\exhibitboxwidth}{L{0.9cm} L{1.2cm} L{1.6cm} X X X}
\toprule
\textbf{代码} & \textbf{名称} & \textbf{主题} & \textbf{订单状态} & \textbf{收入状态} & \textbf{估值影响} \\
\midrule
688146 & 中船特气 & WF6订单、产能、调价机制 & 客户洽谈增加，但长期/大额订单未签署 & 电子特气收入可见，WF6单品拆分未披露 & 可保留最高直接性和事件验证权重，但不能把洽谈或调价机制直接折成增长每股收益。 \\
688549 & 中巨芯-U & WF6产能、扩产、订单和成本 & 无新增长期/大额实质性订单 & WF6收入和毛利未披露 & 直接产能不能升级为订单收入，估值只能保留观察信用和风险折扣。 \\
600378 & 昊华科技 & WF6收入占比和业绩影响 & 订单金额和客户结构未披露 & WF6收入占比0.13\%，单品利润贡献低 & 可保留电子化学品平台价值，但不能给中船特气同等WF6弹性。 \\
688268 & 华特气体 & WF6合成产能和纯化供货 & 按部分客户要求供应，未披露客户名称、金额或长期合同 & 业务占比较小，未披露收入占比 & 电子特气平台可保留，WF6不得写成合成产能或大额订单逻辑。 \\
002971 & 和远气体 & 电子级WF6试生产、订单和业绩 & 未签署实质性订单协议 & 尚未产生任何业绩 & WF6只能作为项目期权和风险提示，不能进入2026E增长每股收益。 \\
\bottomrule
\end{tabularx}
\sourcenote{公司公告PDF、互动平台回复转述和媒体采访归档整理；澄清可以证明未签单、试生产、收入占比低或无合成产能，但不能替代客户侧认证、采购订单、合同平均售价、收入确认比例或单品毛利。}
\end{exhibitbox}

公司澄清层的作用是把市场热度降回可验证字段。中船特气的直接性仍最强，但长期或大额订单尚未公告；中巨芯和和远气体直接挡住了订单和业绩外推；昊华科技把WF6单品贡献限制在很低收入占比；华特气体则明确不是WF6合成产能逻辑。因此，平台需求再强，也只能改变调研优先级和事件跟踪权重，不能越过公司澄清直接进入增长每股收益。

互动问答层进一步补上公司口径的细颗粒度边界。本次巨潮互动易实抓显示：和远气体披露500吨WF6项目已建成并处试生产、尚未形成批量销售；凯美特气、南大光电明确不生产WF6；中钨高新没有进入WF6的计划，但存在向相关客户供应钨粉原料的链条线索；新宙邦的含氟电子材料更多是氟化液和半导体化学品放量线索，不等同WF6订单。上证e互动本轮返回403或非JSON响应，因此中船特气、中巨芯-U、华特气体、昊华科技和金宏气体的互动问答不能计为实抓成功，仍回到公司公告、产品页、年报和风险提示证据层。这些问答能提升传闻否证和后续调研优先级，不能替代合同金额、合同平均售价、收入确认比例、客户份额或单品毛利。

\begin{exhibitbox}[投资者关系互动问答补证（1）]
\scriptsize
\begin{tabularx}{\exhibitboxwidth}{L{0.85cm} L{1.15cm} L{1.55cm} L{1.65cm} X X}
\toprule
\textbf{代码} & \textbf{名称} & \textbf{主题} & \textbf{证据作用} & \textbf{问题摘录} & \textbf{回答摘录} \\
\midrule
002407 & 多氟多 & WF6/NF3；客户/认证；含氟电子材料 & 客户/供货可见度补证 & 贵公司有无建成WF6（六氟化钨）的生产线，未来有无生产六氟化钨的规划？ & 您好，公司现有完整的高纯氟气生产线，具备生产各类与氟气相关特种气体的研发能力和产业化基础，已有高纯氟氮混气、四氟化… \\
002971 & 和远气体 & WF6/NF3；产能/量产 & 负面门禁：否证传闻或阻断业绩外推 & 公司与兴福电子合资的湖北兴远芯气体有限公司目前业务产品有哪些？是否有六氟化钨产能？… & 您好，公司与兴福电子合资的湖北兴远芯气体有限公司主要是做电子大宗气体现场制气业务，不生产六氟化钨产品。公司六氟化钨… \\
300346 & 南大光电 & WF6/NF3；产能/量产 & 产品/产能阶段补证 & 请问贵公司的三氟化氮纯度等级为多少N？总产能如何，是否达产？主要用途是什么，是否可… & 感谢您的关注！公司三氟化氮产品主要应用于IC、LCD制造中的等离子蚀刻和腔体清洗环节，现有产能9200吨。 \\
300037 & 新宙邦 & 订单/合同；价格/调价；客户/认证；产能/量产；含氟电子材料 & 价格边界：仅作合同价格补证线索 & 1.公司的电容器化学品在全国、全球市占率如何？2.公司的半导体湿化产品目前通过了哪… & 在电容化学品方面，公司经过近三十年的深度耕耘，凭借扎实的研发积累、稳定的供应保障、严谨的质量体系等优势，已成为全球… \\
300346 & 南大光电 & WF6/NF3；价格/调价 & 负面门禁：否证传闻或阻断业绩外推 & 董秘你好，六氟化钨最近出口价格涨幅巨大，公司有六氟化钨吗，谢谢 & 感谢您的关注！公司不生产六氟化钨。 \\
002842 & 翔鹭钨业 & 订单/合同；价格/调价；客户/认证；产能/量产；毛利/成本 & 价格边界：仅作合同价格补证线索 & 管理层 好，请问公司光伏金刚线用钨丝业务目前在手订单情况如何？产能利用率大概处于什… & 尊敬的投资者，您好！公司与下游金刚线生产厂家合作关系长期、稳定，在手订单与意向性订单充足。考虑到设备检修、节假日休… \\
\bottomrule
\end{tabularx}
\sourcenote{巨潮互动易和上证e互动公开问答；互动问答只用于公司口径补证、传闻否证、调研优先级和风险折扣，不替代具体订单金额、合同平均售价、收入确认比例、客户份额或单品毛利。完整证据见 data/investor\_relations\_qa\_evidence.json/md。}
\end{exhibitbox}

\begin{exhibitbox}[投资者关系互动问答补证（2）]
\scriptsize
\begin{tabularx}{\exhibitboxwidth}{L{0.85cm} L{1.15cm} L{1.55cm} L{1.65cm} X X}
\toprule
\textbf{代码} & \textbf{名称} & \textbf{主题} & \textbf{证据作用} & \textbf{问题摘录} & \textbf{回答摘录} \\
\midrule
002971 & 和远气体 & WF6/NF3；客户/认证 & 仅提问未回答：不能进入估值 & 尊敬的董秘，您好！贵公司电子特气产品：三氟化氮、六氟化钨、硅烷、乙硅烷、高纯氯化氢… & 未回复 \\
002971 & 和远气体 & WF6/NF3；客户/认证 & 仅提问未回答：不能进入估值 & 尊敬的董秘，您好！贵公司电子特气产品：三氟化氮、六氟化钨、硅烷、乙硅烷、高纯氯化氢… & 未回复 \\
002971 & 和远气体 & WF6/NF3；客户/认证 & 仅提问未回答：不能进入估值 & 尊敬的董秘，您好！贵公司的电子特气产品：三氟化氮、六氟化钨、硅烷、乙硅烷、高纯氯化… & 未回复 \\
002971 & 和远气体 & WF6/NF3；产能/量产 & 负面门禁：否证传闻或阻断业绩外推 & 你好！我司六氟化钨产能现在达到多少了，近期产销情况如何 & 您好，公司该产品尚处于试生产过程，暂未实现稳产量产，离达到正式生产和实现销售还有比较长的过程，请注意风险，谢谢关注！ \\
002407 & 多氟多 & WF6/NF3；客户/认证；产能/量产；含氟电子材料 & 客户/供货可见度补证 & 突发：多氟多午间澄清：公司无六氟化钨（WF₆）（截至2026-06-10），请广大… & 您好，公司现有完整的高纯氟气生产线，具备生产各类与氟气相关特种气体的研发能力和产业化基础，已有高纯氟氮混气、四氟化… \\
002407 & 多氟多 & WF6/NF3；客户/认证；高纯钨材；含氟电子材料 & 客户/供货可见度补证 & 董秘您好，2026年4月28日业绩说明会提到六氟化钨产线已建好、在做客户认证。请问… & 您好，公司现有完整的高纯氟气生产线，具备生产各类与氟气相关特种气体的研发能力和产业化基础，已有高纯氟氮混气、四氟化… \\
\bottomrule
\end{tabularx}
\sourcenote{巨潮互动易和上证e互动公开问答；第二组展示非头部但会影响扩散链判断的电子特气、氟化工和高纯钨材问答。问答证据不替代订单金额、合同平均售价、收入确认比例、客户份额或单品毛利。}
\end{exhibitbox}


投关调研记录层解决的是“这些公司是否存在可追溯的调研纪要、投资者关系活动记录或业绩说明会入口”。本次用巨潮投关专栏和东方财富公告索引做标题级探针，18只标的均有投关、机构调研、业绩说明会或接待日标题命中，其中深市公司多能由巨潮投关专栏返回记录，沪市公司主要由东方财富公告索引补齐标题入口。进一步对东方财富可取正文的记录执行关键词探针，正文覆盖18只标的、69条记录，但P0关闭数量仍为0。这个层级只能证明来源存在、正文语境和后续人工复核优先级；若记录正文没有披露订单金额、合同平均售价、收入确认比例、客户份额或单品毛利，仍不能改变增长业绩和目标价信用。


\begin{exhibitbox}[投关调研记录标题索引与正文关键词探针]
\scriptsize
\begin{tabularx}{\exhibitboxwidth}{L{0.85cm} L{1.15cm} L{1.0cm} L{1.4cm} L{1.35cm} X L{1.5cm}}
\toprule
\textbf{代码} & \textbf{名称} & \textbf{日期} & \textbf{类型} & \textbf{来源} & \textbf{标题} & \textbf{估值处理} \\
\midrule
002407 & 多氟多 & 2026-06-26 & 投资者关系活动记录表 & 巨潮投关专栏 & 2026年6月25日投资者关系活动记录表 & 标题级，不入估值 \\
002407 & 多氟多 & 2026-06-26 & 投资者关系活动记录表 & 东方财富公告索引 & 多氟多:2026年6月25日投资者关系活动记录表 & 标题级，不入估值 \\
600160 & 巨化股份 & 2026-06-18 & 投资者网上接待日 & 东方财富公告索引 & 巨化股份:巨化股份关于举办投资者接待日活动的公告 & 标题级，不入估值 \\
688268 & 华特气体 & 2026-06-15 & 投资者关系活动记录表 & 东方财富公告索引 & 华特气体:2026-004广东华特气体股份有限公司投资者关系活动记录表(2026年6月12日) & 标题级，不入估值 \\
603505 & 金石资源 & 2026-06-13 & 业绩说明会 & 东方财富公告索引 & 金石资源:金石资源集团股份有限公司关于2025年度暨2026年第一季度沪市主板生态提效能之稳利筑基专题集体… & 标题级，不入估值 \\
600160 & 巨化股份 & 2026-06-11 & 投资者关系活动记录表 & 东方财富公告索引 & 巨化股份:投资者关系活动记录表(2026年6月9日) & 标题级，不入估值 \\
603505 & 金石资源 & 2026-06-06 & 业绩说明会 & 东方财富公告索引 & 金石资源:金石资源集团股份有限公司关于参加2025年度暨2026年第一季度沪市主板生态提效能之稳利筑基专题… & 标题级，不入估值 \\
688268 & 华特气体 & 2026-05-26 & 投资者关系活动记录表 & 东方财富公告索引 & 华特气体:2026-003广东华特气体股份有限公司投资者关系活动记录表(2026年5月26日) & 标题级，不入估值 \\
688146 & 中船特气 & 2026-05-26 & 业绩说明会 & 东方财富公告索引 & 中船特气:中船特气关于召开2025年年度暨2026年第一季度业绩说明会的公告 & 标题级，不入估值 \\
688106 & 金宏气体 & 2026-05-22 & 投资者关系活动记录表 & 东方财富公告索引 & 金宏气体:投资者关系活动记录表2026年5月20日 & 标题级，不入估值 \\
605020 & 永和股份 & 2026-05-22 & 业绩说明会 & 东方财富公告索引 & 永和股份:浙江永和制冷股份有限公司关于召开2025年度暨2026年第一季度业绩说明会的公告 & 标题级，不入估值 \\
603379 & 三美股份 & 2026-05-22 & 业绩说明会 & 东方财富公告索引 & 三美股份:浙江三美化工股份有限公司关于召开2025年度暨2026年第一季度业绩说明会的公告 & 标题级，不入估值 \\
\bottomrule
\end{tabularx}
\sourcenote{巨潮投关专栏、东方财富公告索引和东方财富公告正文接口；标题/索引和正文关键词探针只证明来源存在、正文语境和人工复核优先级，不替代订单金额、合同平均售价、收入确认比例、客户份额或单品毛利。完整证据见 data/ir\_activity\_record\_evidence.json/md。}
\end{exhibitbox}


问询/审核回复层比投关记录更正式，也比普通公告标题更接近硬字段复核。本轮扫描交易所问询、审核问询回复、反馈意见、监管函、关注函和落实函，覆盖18只标的，命中3只、解析13条。它的用途不是把“问询过”写成“订单已经兑现”，而是确认监管或交易所文件是否披露了产品收入、客户订单、合同平均售价、收入确认、单品毛利或募投产能等可量化字段。若正文只有募投说明、泛化业务描述或未披露硬字段，仍不改变增长每股收益信用。


\begin{exhibitbox}[问询/审核回复公告证据]
\scriptsize
\begin{tabularx}{\exhibitboxwidth}{L{0.85cm} L{1.15cm} L{1.0cm} L{1.25cm} X X X}
\toprule
\textbf{代码} & \textbf{名称} & \textbf{日期} & \textbf{类型} & \textbf{标题} & \textbf{硬字段状态} & \textbf{估值处理} \\
\midrule
002842 & 翔鹭钨业 & 2026-06-16 & 审核问询回复 & 翔鹭钨业:发行人及保荐人关于审核问询函的回复 & 正文命中客户/订单、价格/毛利、确认/回款、产能/募投语境，可作为硬缺口复核入口 & 硬字段复核；不直接入模 \\
002842 & 翔鹭钨业 & 2026-06-16 & 审核问询回复 & 翔鹭钨业:广东翔鹭钨业股份有限公司关于向特定对象发行股票审核问询函回复及募集说明书等申请文件更… & 标题命中问询/回复，但正文未命中核心硬字段语境 & 硬字段复核；不直接入模 \\
002842 & 翔鹭钨业 & 2026-06-15 & 审核问询回复 & 7-1发行人及保荐人关于审核问询函的回复(广东翔鹭钨业股份有限公司) & 正文命中客户/订单、价格/毛利、确认/回款、产能/募投语境，可作为硬缺口复核入口 & 硬字段复核；不直接入模 \\
002842 & 翔鹭钨业 & 2025-12-03 & 审核问询回复 & 翔鹭钨业:会计师关于审核问询函的回复 & 标题命中问询/回复，但正文未命中核心硬字段语境 & 硬字段复核；不直接入模 \\
002842 & 翔鹭钨业 & 2025-12-03 & 审核问询回复 & 翔鹭钨业:发行人及保荐人关于审核问询函的回复 & 正文命中客户/订单、价格/毛利、确认/回款、产能/募投语境，可作为硬缺口复核入口 & 硬字段复核；不直接入模 \\
002971 & 和远气体 & 2025-08-21 & 审核问询回复 & 和远气体:关于向特定对象发行股票审核问询函回复等文件更新的提示性公告 & 标题命中问询/回复，但正文未命中核心硬字段语境 & 硬字段复核；不直接入模 \\
002971 & 和远气体 & 2025-08-21 & 审核问询回复 & 和远气体:立信会计师事务所(特殊普通合伙)关于湖北和远气体股份有限公司申请向特定对象发行股票的… & 正文命中价格/毛利、确认/回款语境，可作为硬缺口复核入口 & 硬字段复核；不直接入模 \\
002971 & 和远气体 & 2025-08-21 & 审核问询回复 & 和远气体:关于湖北和远气体股份有限公司申请向特定对象发行股票的审核问询函回复 & 正文命中确认/回款语境，可作为硬缺口复核入口 & 硬字段复核；不直接入模 \\
\bottomrule
\end{tabularx}
\sourcenote{东方财富公告索引与正文接口；扫描问询函、审核问询回复、反馈意见、监管函、关注函和落实函。问询/审核回复证据只提供交易所或监管层硬字段复核入口，不得把问询标题、募投说明或泛化回复替代单品订单、合同平均售价、收入确认比例、客户份额、单品收入或单品毛利。完整证据见 data/regulatory\_inquiry\_evidence.json/md。}
\end{exhibitbox}


合同公告层比投关记录更接近P0字段，但本轮扫描结果仍然支持“部分补证、不能直接上修”。2024年以来18只标的的重大合同、中标、供货、长单、框架协议和合作协议公告中，命中9只标的、解析25条正文，其中WF6/NF3主题单品订单金额样本2条，电子气体平台订单金额样本1条。最有信息量的是中船特气：2025年披露境内某知名集成电路客户A六氟化钨采购合同约1.1904亿元（不含税），2024年披露三氟化氮销售合同约1.1508亿元（含税），2025年还披露HF/HCL/D2等11种产品采购合同约1.4822亿元（不含税）。前两项能把中船特气的历史WF6/NF3订单金额从“未找到”升级为“部分补证”，第三项说明电子气体平台订单可见度更强；但三者仍不能替代合同平均售价、收入确认比例、客户份额和单品毛利，也不能证明2026年以后订单持续放大。和远气体、金宏气体的大额合同主要是氧气、氮气、丙烷等工业现场制气，新宙邦是锂电电解液协议，厦门钨业多为钼精矿/稀土关联交易框架，这些都不能迁移到电子特气增长每股收益。


\begin{exhibitbox}[合同/中标/供货协议公告全文扫描]
\scriptsize
\begin{tabularx}{\exhibitboxwidth}{L{0.85cm} L{1.15cm} L{1.0cm} L{1.25cm} X X X}
\toprule
\textbf{代码} & \textbf{名称} & \textbf{日期} & \textbf{类型} & \textbf{标题} & \textbf{产品/主题范围} & \textbf{P0状态} \\
\midrule
688146 & 中船特气 & 2024-07-18 & 合同公告 & 中船特气:中船特气关于自愿披露签订日常重要合同的公告 & WF6/NF3正式合同，订单金额已披露但平均售价/毛利未披露 & 披露WF6/NF3合同金额和匿名集成电路客户；订单金额部分补证，但平均售价/确认比… \\
688146 & 中船特气 & 2025-04-09 & 合同公告 & 中船特气:中船特气关于自愿披露获得重要采购合同的公告 & WF6/NF3正式合同，订单金额已披露但平均售价/毛利未披露 & 披露WF6/NF3合同金额和匿名集成电路客户；订单金额部分补证，但平均售价/确认比… \\
688146 & 中船特气 & 2025-04-23 & 合同公告 & 中船特气:中船特气关于自愿披露签订日常重要合同的公告 & 电子气体平台合同，披露金额但不等同WF6/NF3单品订单 & 披露电子气体平台合同金额和匿名集成电路客户；补强平台订单可见度，但不关闭WF6/N… \\
002407 & 多氟多 & 2024-06-18 & 合作/框架协议 & 多氟多:关于签署框架合作协议的提示性公告 & 主题相关，需要人工复核是否为电子特气/电子材料单品 & 未发现可直接入模的合同金额、单品价格、收入确认比例或毛利字段 \\
002407 & 多氟多 & 2024-09-13 & 合作/框架协议 & 多氟多:关于签署框架合作协议的进展公告 & 合同/订单公告，主题相关性需人工复核 & 未发现可直接入模的合同金额、单品价格、收入确认比例或毛利字段 \\
002407 & 多氟多 & 2025-11-07 & 订单/长单线索 & 多氟多:关于公司产品取得客户订单的公告 & 合同/订单公告，主题相关性需人工复核 & 未发现可直接入模的合同金额、单品价格、收入确认比例或毛利字段 \\
002971 & 和远气体 & 2024-08-23 & 合作/框架协议 & 和远气体:关于公司与武汉东湖综合保税区建设服务中心签订投资合作协议的公告 & 主题相关，需要人工复核是否为电子特气/电子材料单品 & 未发现可直接入模的合同金额、单品价格、收入确认比例或毛利字段 \\
002971 & 和远气体 & 2025-04-22 & 供气合同 & 和远气体:关于公司与宜昌邦普时代新能源有限公司签订供气合同的公告 & 合同/订单公告，主题相关性需人工复核 & 披露定价/期限/订单框架线索，但未给可入模单品金额、均价或毛利 \\
\bottomrule
\end{tabularx}
\sourcenote{东方财富公告索引和公告正文接口；合同公告扫描用于审计官方公告层是否存在金额、期限、定价或产品范围线索。中船特气历史WF6/NF3合同补强主题单品订单金额样本，HF/HCL/D2等11种产品合同补强电子气体平台订单可见度，但均不替代合同平均售价、收入确认比例、客户份额或单品毛利；工业气体、原料采购、关联交易和锂电协议不得替代WF6/NF3/电子材料单品订单。完整证据见 data/contract\_order\_announcement\_evidence.json/md。}
\end{exhibitbox}

合同金额经济性约束把这些金额放到公司收入分母里看。中船特气2025年电子特种气体分部收入为19.27亿元，公司总收入约22.60亿元；2025年WF6合同约1.1904亿元、2024年NF3合同约1.1508亿元、2025年HF/HCL/D2等平台合同约1.4822亿元，单笔金额相对公司收入和电子特气分部收入均有可见强度。但这个比例桥只说明历史订单金额的量级和事件验证强度，不能解决平均售价、销量拆分、交付节奏、收入确认比例、客户份额和单品毛利，所以仍不能上修长期增长每股收益。

\begin{exhibitbox}[合同金额经济性约束]
\scriptsize
\begin{tabularx}{\exhibitboxwidth}{L{1.0cm} L{1.7cm} L{1.15cm} L{0.8cm} L{1.0cm} L{1.15cm} X X}
\toprule
\textbf{日期} & \textbf{产品范围} & \textbf{金额} & \textbf{税口径} & \textbf{金额/收入} & \textbf{金额/电子特气} & \textbf{证据等级} & \textbf{仍阻断字段} \\
\midrule
2025-04-23 & HF/HCL/D2等电子气体平台 & 1.5亿元 & 不含税 & 6.6\% & 7.7\% & 电子气体平台订单金额补证 & 合同平均售价、销量/单价拆分、交付节奏、收入确认比例、客户份额、单品毛利 \\
2025-04-09 & 六氟化钨/WF6 & 1.2亿元 & 不含税 & 5.3\% & 6.2\% & 主题单品订单金额部分补证 & 合同平均售价、销量/单价拆分、交付节奏、收入确认比例、客户份额、单品毛利 \\
2024-07-18 & 三氟化氮/NF3 & 1.2亿元 & 含税 & 5.1\% & 6.0\% & 主题单品订单金额部分补证 & 合同平均售价、销量/单价拆分、交付节奏、收入确认比例、客户份额、单品毛利 \\
\bottomrule
\end{tabularx}
\sourcenote{合同金额来自中船特气自愿披露合同公告，收入分母来自2025年报分产品表和主营构成。含税金额与收入口径不可直接同口径比较；本表只约束历史订单强度和收入桥上限，不替代合同平均售价、收入确认比例、客户份额或单品毛利。完整证据见 data/contract\_economics\_constraint.json/md。}
\end{exhibitbox}

订单耐久性门禁把“历史发生过订单”和“可以进入持续增长每股收益”分开。中船特气的WF6/NF3历史合同金额样本和电子气体平台订单，确实把订单证据从传闻或洽谈推进到公告金额层；但耐久性门禁仍显示18只标的中没有一家能直接进入持续订单增长每股收益。原因是持续订单、合同平均售价、交付节奏、收入确认比例、客户份额、单品收入和单品毛利没有同时披露。非主题工业气体合同、原料框架、锂电协议、合同负债和预约披露日期，只能作为来源穷尽或验证窗口，不能替代电子特气或电子材料单品订单。

\begin{exhibitbox}[订单持续性与收入确认耐久性证据]
\scriptsize
\begin{tabularx}{\exhibitboxwidth}{L{0.85cm} L{1.15cm} L{1.45cm} X L{1.65cm} L{1.5cm} X}
\toprule
\textbf{代码} & \textbf{名称} & \textbf{订单阶段} & \textbf{订单证据} & \textbf{耐久性门禁} & \textbf{验证窗口} & \textbf{仍缺字段} \\
\midrule
688146 & 中船特气 & 历史主题订单金额样本 & 六氟化钨/WF61.2亿元；三氟化氮/NF31.2亿元；HF/HCL/D2等电子气体平台1.5亿元 & 事件验证；持续性未闭合 & 2026年 半年报 2026-07-18 & 持续订单、合同平均售价、收入确认比例、客户份额、单品收入和单品毛利。 \\
002971 & 和远气体 & 公司澄清负面边界 & 未签署实质性订单协议 & 负面边界；订单阻断 & 2026年 一季报 2026-04-24 & 法律约束订单、批量销售、合同平均售价、收入确认比例和单品毛利。 \\
688549 & 中巨芯-U & 公司澄清负面边界 & 无新增长期/大额实质性订单 & 负面边界；订单阻断 & 2026年 半年报 2026-08-08 & 法律约束订单、批量销售、合同平均售价、收入确认比例和单品毛利。 \\
002407 & 多氟多 & 非主题/弱相关合同公告 & 合同/中标/供货协议公告解析3条，但产品范围不能迁移为WF6/NF3或电子材料单品订单 & 非主题合同；不得迁移 & 2026年 一季报 2026-04-24 & 主题单品订单、合同价格、客户份额、收入确认比例和单品毛利。 \\
300037 & 新宙邦 & 非主题/弱相关合同公告 & 合同/中标/供货协议公告解析2条，但产品范围不能迁移为WF6/NF3或电子材料单品订单 & 非主题合同；不得迁移 & 2026年 一季报 2026-04-28 & 主题单品订单、合同价格、客户份额、收入确认比例和单品毛利。 \\
600160 & 巨化股份 & 非主题/弱相关合同公告 & 合同/中标/供货协议公告解析2条，但产品范围不能迁移为WF6/NF3或电子材料单品订单 & 非主题合同；不得迁移 & 2026年 半年报 2026-08-26 & 主题单品订单、合同价格、客户份额、收入确认比例和单品毛利。 \\
600378 & 昊华科技 & 非主题/弱相关合同公告 & 合同/中标/供货协议公告解析1条，但产品范围不能迁移为WF6/NF3或电子材料单品订单 & 非主题合同；不得迁移 & 2026年 半年报 2026-08-27 & 主题单品订单、合同价格、客户份额、收入确认比例和单品毛利。 \\
600549 & 厦门钨业 & 非主题/弱相关合同公告 & 合同/中标/供货协议公告解析5条，但产品范围不能迁移为WF6/NF3或电子材料单品订单 & 非主题合同；不得迁移 & 2026年 半年报 2026-08-21 & 主题单品订单、合同价格、客户份额、收入确认比例和单品毛利。 \\
688106 & 金宏气体 & 非主题/弱相关合同公告 & 合同/中标/供货协议公告解析4条，但产品范围不能迁移为WF6/NF3或电子材料单品订单 & 非主题合同；不得迁移 & 2026年 半年报 2026-08-26 & 主题单品订单、合同价格、客户份额、收入确认比例和单品毛利。 \\
300346 & 南大光电 & 收入确认代理存在 & 合同负债/收入2.6\%；确认门槛：签收/领用清单/报关提单 & 确认代理；订单未闭合 & 2026年 一季报 2026-04-24 & 单品订单金额、合同价格、交付节奏、客户份额和单品毛利。 \\
688268 & 华特气体 & 收入确认代理存在 & 合同负债/收入7.8\%；确认门槛：验收签收/寄售对账/报关离岸 & 确认代理；订单未闭合 & 2026年 半年报 2026-08-26 & 单品订单金额、合同价格、交付节奏、客户份额和单品毛利。 \\
000657 & 中钨高新 & 未找到主题订单 & 公开语料内未找到主题单品合同、长单、中标或持续供货金额证据 & 未找到主题订单 & 2026年 一季报 2026-04-27 & 主题单品订单、合同价格、收入确认比例、客户份额和单品毛利。 \\
002378 & 章源钨业 & 未找到主题订单 & 公开语料内未找到主题单品合同、长单、中标或持续供货金额证据 & 未找到主题订单 & 2026年 一季报 2026-04-28 & 主题单品订单、合同价格、收入确认比例、客户份额和单品毛利。 \\
002549 & 凯美特气 & 未找到主题订单 & 公开语料内未找到主题单品合同、长单、中标或持续供货金额证据 & 未找到主题订单 & 2026年 一季报 2026-04-22 & 主题单品订单、合同价格、收入确认比例、客户份额和单品毛利。 \\
002842 & 翔鹭钨业 & 未找到主题订单 & 公开语料内未找到主题单品合同、长单、中标或持续供货金额证据 & 未找到主题订单 & 2026年 一季报 2026-04-29 & 主题单品订单、合同价格、收入确认比例、客户份额和单品毛利。 \\
603379 & 三美股份 & 未找到主题订单 & 公开语料内未找到主题单品合同、长单、中标或持续供货金额证据 & 未找到主题订单 & 2026年 半年报 2026-08-25 & 主题单品订单、合同价格、收入确认比例、客户份额和单品毛利。 \\
603505 & 金石资源 & 未找到主题订单 & 公开语料内未找到主题单品合同、长单、中标或持续供货金额证据 & 未找到主题订单 & 2026年 半年报 2026-08-21 & 主题单品订单、合同价格、收入确认比例、客户份额和单品毛利。 \\
605020 & 永和股份 & 未找到主题订单 & 公开语料内未找到主题单品合同、长单、中标或持续供货金额证据 & 未找到主题订单 & 2026年 半年报 2026-08-17 & 主题单品订单、合同价格、收入确认比例、客户份额和单品毛利。 \\
\bottomrule
\end{tabularx}
\sourcenote{订单耐久性门禁由合同公告扫描、合同金额经济性约束、官方单品收入/订单金额边界、公司澄清、收入确认代理和业绩披露日历合成；该表只区分历史订单、平台订单、非主题合同、负面边界和确认代理，不替代持续订单、合同平均售价、收入确认比例、客户份额、单品收入或单品毛利。完整证据见 data/order\_durability\_evidence.json/md。}
\end{exhibitbox}

合同公告、公司澄清和招股书证据合并后，单品层面的结论更清楚：中船特气已经不是“完全没有订单金额证据”，而是“有历史WF6/NF3金额样本、但缺平均售价和毛利”；昊华科技已经不是“只有产品页”，而是“2025年WF6收入占比只有0.13\%”；中巨芯-U、华特气体和和远气体则分别受无新增长期大额订单、无WF6合成产能、尚未产生业绩和无实质性订单约束。因此本报告将电子特气标的继续放入事件验证池，而不是直接把产业链热度写成基本面目标价上修。

\begin{exhibitbox}[官方单品收入/订单金额边界]
\scriptsize
\begin{tabularx}{\exhibitboxwidth}{L{0.85cm} L{1.15cm} L{1.45cm} X X X X}
\toprule
\textbf{代码} & \textbf{名称} & \textbf{单品/范围} & \textbf{官方量化字段} & \textbf{可证明} & \textbf{P0状态} & \textbf{仍缺字段} \\
\midrule
600378 & 昊华科技 & WF6产品收入占比 & 收入占比0.13\% & WF6产品真实存在，但对2025年整体收入贡献极低。 & 部分补证：收入占比可见，盈利字段未关闭 & 单品订单、合同平均售价、客户份额、收入确认比例、单品毛利 \\
688146 & 中船特气 & WF6主题单品合同 & 1.1904亿元 & 官方公告层存在WF6主题单品订单金额样本。 & 部分补证：订单金额样本可见，盈利字段未关闭 & 合同平均售价、交付节奏、收入确认比例、客户份额、单品毛利 \\
688146 & 中船特气 & NF3主题单品合同 & 1.1508亿元 & 官方公告层存在NF3主题单品订单金额样本，能补强电子特气订单可见度。 & 部分补证：订单金额样本可见，盈利字段未关闭 & 合同平均售价、交付节奏、收入确认比例、客户份额、单品毛利 \\
688146 & 中船特气 & 电子气体平台合同 & 1.4822亿元 & 电子气体平台订单可见度提高，但产品范围不是WF6/NF3单品。 & 平台补证：不关闭WF6/NF3单品P0 & WF6/NF3单品订单、合同平均售价、收入确认比例、客户份额、单品毛利 \\
688146 & 中船特气 & 历史WF6产销 & 历史销量约269-838吨 & WF6历史产销兑现和产能利用可被一级来源量化。 & 历史补证：不等于当前收入拆分 & 当前订单、合同平均售价、收入确认比例、客户份额、当前单品毛利 \\
688146 & 中船特气 & 价格/毛利披露边界 & 价格降幅8.27\%/5.76\%/5.16\%；毛利率下降7.50pct & 公开资料能证明价格和毛利方向，但不能得到绝对平均售价或当前单品毛利。 & 披露边界：平均售价和毛利仍阻断 & 当前绝对平均售价、当前单品毛利、价格传导条款 \\
688549 & 中巨芯-U & WF6产能/订单边界 & 产能600吨；订单未关闭 & 具备直接产品能力，但订单和盈利外推被公司口径限制。 & 负面边界：收入/订单未披露 & 订单金额、合同平均售价、客户份额、收入确认比例、单品收入、单品毛利 \\
688268 & 华特气体 & WF6纯化供货边界 & 未披露占比 & 公司不应按WF6合成产能逻辑估值，纯化供货弹性有限。 & 负面边界：非合成产能逻辑 & 纯化供货收入占比、订单金额、合同平均售价、客户份额、单品毛利 \\
002971 & 和远气体 & WF6/NF3试生产边界 & 0业绩贡献/无实质性订单 & 试生产和项目存在不能转成收入或订单。 & 负面边界：业绩与订单均未关闭 & 批量销售、法律约束订单、合同平均售价、收入确认、单品毛利 \\
\bottomrule
\end{tabularx}
\sourcenote{公司澄清、东方财富公告正文扫描和中船特气招股说明书归档整理；该表只收窄单品收入/订单金额边界，不替代合同平均售价、收入确认比例、客户份额或单品毛利。完整证据见 data/official\_single\_product\_revenue\_boundary.json/md。}
\end{exhibitbox}

\begin{exhibitbox}[客户链高影响声明审计]
\footnotesize
\begin{tabularx}{\exhibitboxwidth}{L{1.3cm} X L{1.35cm} L{0.85cm} L{0.95cm} X}
\toprule
\textbf{公司} & \textbf{声明} & \textbf{来源} & \textbf{置信度} & \textbf{估值采用} & \textbf{正文表述} \\
\midrule
主题层 & WF6用于半导体CVD钨沉积，是钨粉、氟化工和电子特气交汇点 & S2 & 高 & 是 & 可作为产业链交汇点证据，不能单独证明任何公司的客户或订单 \\
中船特气 & 中船特气WF6现有产能2000吨/年、6N级 & S6, S9, S18, S32, S33 & 高 & 是 & 产能和规格直接性高；历史WF6/NF3合同金额样本已披露（2025-04-09六氟化钨合同约1.1904亿元；2024-07-18三氟化氮合同约1.1508亿元），但合同平均售价、收入确认比例、客户份额、单品毛利和2026年后持续订单未披露 \\
中船特气 & 中船特气披露历史WF6/NF3主题单品合同金额样本 & S32, S33 & 高 & 是 & 订单金额可从纯缺口升级为历史部分补证：2025-04-09六氟化钨合同约1.1904亿元；2024-07-18三氟化氮合同约1.1508亿元；但不得替代合同平均售价、收入确认比例、客户份额或单品毛利，也不得直接进入长期增长每股收益 \\
中船特气 & 中船特气披露电子气体平台合同金额样本 & S34 & 高 & 是 & 2025-04-23披露HF/HCL/D2等11种电子气体产品采购合同约1.4822亿元；该证据补强电子气体平台订单可见度，但不得替代WF6/NF3单品订单、合同平均售价、收入确认比例、客户份额或单品毛利，也不得直接进入长期增长每股收益 \\
中船特气 & 中船特气3383吨电子气体项目有公告和环评公示证据 & S29, S30, S31 & 高 & 是 & 项目、投资额、分期总产能和六氟化钨六氟化钼车间可进入扩产监测；不得替代WF6单品产能拆分、订单、合同平均售价或单品毛利 \\
中巨芯-U & 中巨芯高纯WF6产能600吨，暂无扩产计划 & S10, S14 & 高 & 是 & 直接产能存在，但长期/大额订单和扩产边界限制估值信用 \\
昊华科技 & 昊华科技2025年度WF6收入占总收入0.13\% & S2, S11, S15 & 高 & 是 & 有WF6产品能力，但收入占比低，不能给中船同等弹性 \\
华特气体 & 华特气体无WF6合成产能，仅按客户要求供应纯化WF6，业务占比较小 & S12, S17 & 高 & 是 & 保留电子特气平台价值，但不得把WF6合成产能或大额订单假设写入估值 \\
和远气体 & 和远气体WF6尚处试生产，尚未产生任何业绩且未签署实质性订单 & S16 & 高 & 是 & 只可作为项目期权和风险提示，不得进入2026E增长每股收益 \\
主题层 & 2026年4月中国WF6出口均价149.79美元/千克，同比+28.33\%、环比+203.83\%，可作为主题热度证据但不能替代公司合同平均售价 & S13 & 中 & 是 & 行业价格进入情绪锚和监测变量，不直接填入单公司收入桥或增长每股收益 \\
华特气体 & 华特气体是电子特气平台并有历史券商目标价 & S8 & 中 & 是 & 可作为电子特气平台证据；WF6直接收入和客户需另证 \\
电子气体扩散池 & 所有电子气体公司均直接受益WF6涨价 & 未找到 & 低 & 否 & 禁止进入正文强结论；必须逐家公司看产品、收入、客户和订单 \\
\bottomrule
\end{tabularx}
\sourcenote{本表展示影响估值信用的核心客户链声明；年报客户认证/具名客户披露明细另见第5章“客户认证/具名客户披露证据”，完整20条审计记录见 data/customer\_chain\_audit.json/md。}
\end{exhibitbox}

\section{三段利润传导的差异}
钨资源段的利润传导最直接：价格上涨先进入矿山、APT和钨粉环节，再向硬质合金和刀具扩散。问题在于钨价越高，下游补库越谨慎，加工企业可能出现收入增长但库存和现金流承压。因此钨股不能只看价格弹性，还要看资源自给率、库存成本、长单报价和产品结构。

氟化工段有两条完全不同的利润线。制冷剂是配额约束下的供给侧景气，盈利兑现比半导体材料更直接；电子级HF、氟化液和含氟电子特气则需要客户认证和良率验证，兑现周期更长。生态环境部2026年度HFC配额附件已经把巨化、三美、永和体系的制冷剂配额补到公司/实体层级，因此这三家的短期利润更应由R32/R125/R134a价差和配额执行解释；新宙邦和多氟多则需要拆分有机氟、电子化学品和锂电材料周期，不能套用制冷剂配额龙头逻辑。

电子特气段的估值弹性最高，但最怕“名义产能”。真正有价值的是可稳定供应的高纯品级、容器和纯化控制、晶圆厂认证、长协订单、调价机制和客户结构。中船特气因公司公告披露WF6现有产能2000吨/年、6N级而具备最强直接性；中巨芯具备高纯WF6产能，但公告同时提示无新增长期/大额订单和暂无扩产计划；昊华科技有WF6产品页和平台能力，但公告披露2025年WF6收入占比低，因此不能把中船的弹性直接套给昊华。

\section{标的分组打法}
组合研究上可以分三组跟踪。第一组是业绩兑现组：厦门钨业、章源钨业、巨化股份、三美股份，重点看二/三季度利润和价格高位持续性。第二组是事件验证组：中船特气、中巨芯、华特气体、南大光电，重点看公告、调研纪要、客户认证和长单。第三组是扩散观察组：金宏气体、凯美特气、和远气体、多氟多、新宙邦、金石资源，重点看是否从主题扩散走向具体产品收入。
